\documentclass[11pt,a4paper]{article}
\usepackage[T1]{fontenc}
\usepackage[utf8]{inputenc}
\usepackage{lmodern}
\usepackage{amsmath,amssymb,amsthm}
\usepackage{geometry}
\usepackage{microtype}
\usepackage{hyperref}
\geometry{left=23mm,right=23mm,top=23mm,bottom=23mm,headheight=14pt}
\hypersetup{hidelinks,pdftitle={The Coordinate Pythagorean Theorem from Planar Isotropy: A Lie-Theoretic and Axiomatic Derivation},pdfauthor={ChatGPT (GPT-6, High thinking effort); prompted by Alessandro Odasso},pdfsubject={AI-generated mathematical manuscript with authorship disclosure}}
\newtheorem{lemma}{Lemma}[section]
\newtheorem{theorem}{Theorem}[section]
\newtheorem{remark}{Remark}[section]
\newtheorem{axiom}{Axiom}[section]
\renewcommand{\proofname}{Proof}
\title{\textbf{The Coordinate Pythagorean Theorem from Planar Isotropy:}\\[0.25em]
\large A Lie-Theoretic and Axiomatic Derivation}
\author{\normalsize AI-generated manuscript: \textbf{ChatGPT (GPT-6, High thinking effort)}\\
\normalsize Research question and initiating prompt: \textbf{Alessandro Odasso}}
\date{Original draft: 18 August 2026\\Disclosure revision: 12 October 2026}
\begin{document}
\maketitle
\begin{abstract}
We present an axiomatic proof of the Pythagorean theorem in coordinate space derived from primitive Euclidean congruence and continuous rotational symmetry. By representing continuous planar rotations by the one-parameter flow group $G=\{e^{tJ}:t\in\mathbb{R}\}$ generated by
$J=\begin{pmatrix}0&-1\\1&0\end{pmatrix}$, we show that a nonzero symmetric quadratic form $Q$ is invariant under $G$ if and only if $Q(x,y)=\lambda(x^2+y^2)$ ($\lambda\ne0$). By establishing two logically independent channels---a dynamical invariant channel from $G$-flow and a geometric congruence channel from rigid rotation and coordinate calibration---we prove $a^2+b^2=c^2$ without relying on area dissections, trigonometric identities, or an a priori Cartesian distance formula. The manuscript was generated by ChatGPT in response to a research prompt, as disclosed in Section~\ref{sec:introduction}; a request to seek an original proof is not a verified claim of historical novelty.
\end{abstract}

\section{Introduction}\label{sec:introduction}
\noindent\textbf{Origin, AI authorship, and research prompt.}
This manuscript was written by OpenAI's ChatGPT, using the \textbf{GPT-6 model with High thinking effort}, following a prompt from \textbf{Alessandro Odasso} asking whether ChatGPT could find an \emph{original proof of the Pythagorean theorem}. Odasso supplied the initiating research question; the prose and proposed mathematical derivation in the original draft were generated by AI. The word \emph{original} describes the objective of the prompt, \emph{not} an independently established claim of novelty, priority, or publication-level originality. The mathematical argument should be evaluated on its explicit assumptions and deductions, independently of the system that generated the text. This revision adds provenance disclosure while preserving the original argument's conditional axiomatic character.

Many familiar proofs of the Pythagorean theorem use geometric dissection, similarity, or trigonometric/algebraic identities. When extending the theorem to analytic geometry, standard presentations frequently define the Cartesian metric distance via $d(P_1,P_2)=\sqrt{\Delta x^2+\Delta y^2}$, introducing potential circularity if used to evaluate diagonal segment lengths.

In this paper, we demonstrate that the coordinate expression $a^2+b^2=c^2$ arises from planar isotropy together with primitive rigid congruence and coordinate calibration. Rotational symmetry generated by $J$ forces the quadratic form $x^2+y^2$ up to scale, while primitive rigid congruence independently identifies the terminal axial coordinate of a rotated diagonal vector with its geometric length.

\section{Axiomatic Foundations}
We formalize the relationship between primitive Euclidean geometry and Cartesian coordinates through two axioms.

\begin{axiom}[Orthogonal Coordinate Calibration]
The two coordinate directions correspond to perpendicular geometric directions. Furthermore, any two points in $\mathbb{R}^2$ whose coordinate difference vector is $(L,0)$ or $(0,L)$ have a primitive Euclidean geometric separation equal to $|L|$.
\end{axiom}

\begin{axiom}[Rigid Rotational Congruence and Transitivity]
Continuous rotation about the origin $C(0,0)$ is a length-preserving rigid motion belonging to the one-parameter flow group
\begin{equation}\label{eq:flow}
G=\{R(t)=e^{tJ}:t\in\mathbb{R}\},\qquad J=\begin{pmatrix}0&-1\\1&0\end{pmatrix}.
\end{equation}
Furthermore, for any two rays emanating from $C$, there exists a transformation $R(T)\in G$ mapping one ray onto the other.
\end{axiom}

\section{Classification of Invariant Quadratic Forms under \texorpdfstring{$G$}{G}}
We establish that the infinitesimal generator $J$ forces the form of the Euclidean quadratic invariant up to an overall scalar multiplier.

\begin{lemma}[Classification of Nonzero Symmetric Quadratic Invariants]
Let $Q(\mathbf r)=\mathbf r^{T}M\mathbf r$ be a nonzero symmetric quadratic form generated by $M=\begin{pmatrix}\alpha&\beta\\\beta&\gamma\end{pmatrix}$. Then $Q$ is invariant under the rotational flow $G$ if and only if $M=\lambda I$ for some scalar $\lambda\ne0$.
\end{lemma}
\begin{proof}
By definition, $Q$ is invariant under $R(t)=e^{tJ}$ if and only if $R(t)^TMR(t)=M$ for all $t\in\mathbb{R}$. Differentiating at $t=0$ yields the infinitesimal condition
\begin{equation}\label{eq:infinitesimal}J^TM+MJ=0.\end{equation}
\textbf{Uniqueness.} Since $J^T=-J$, equation~\eqref{eq:infinitesimal} reduces to $MJ=JM$. Direct multiplication gives
\begin{equation}\label{eq:products}
MJ=\begin{pmatrix}\beta&-\alpha\\\gamma&-\beta\end{pmatrix},\qquad
JM=\begin{pmatrix}-\beta&-\gamma\\\alpha&\beta\end{pmatrix}.
\end{equation}
Equating corresponding entries forces $\beta=0$ and $\alpha=\gamma$. Setting $\alpha=\gamma=\lambda$ gives $M=\lambda I$. Since $Q\ne0$, we have $\lambda\ne0$.

\textbf{Existence.} Conversely, suppose $M=\lambda I$ with $\lambda\ne0$. Then
\begin{equation}
J^TM+MJ=\lambda(J^T+J)=\lambda(-J+J)=0.
\end{equation}
Differentiating $Q(\mathbf r(t))$ along any integral curve $\dot{\mathbf r}=J\mathbf r$ yields
\begin{equation}
\frac{d}{dt}Q(\mathbf r(t))=\mathbf r^T(J^TM+MJ)\mathbf r=0.
\end{equation}
Thus every quadratic form $Q(x,y)=\lambda(x^2+y^2)$ with $\lambda\ne0$ is conserved along the flow $G$.
\end{proof}

\begin{remark}
The identities $R(t)^TR(t)=I$ and $\det R(t)=1$ follow from $J^T=-J$ and $\operatorname{tr}J=0$. Thus $G$ is a one-parameter subgroup of $SO(2)$ generated by $J$.
\end{remark}

\section{Main Result: The Dual-Channel Proof}
\begin{theorem}[Coordinate Pythagorean Theorem]
In a calibrated orthogonal coordinate frame with primitive length congruence, the length-preserving flow $e^{tJ}$, and ray transitivity, let $\triangle CAB$ be a right triangle with leg lengths $|CA|=a$ and $|AB|=b$, and hypotenuse length $|CB|=c$. Then
\[
a^2+b^2=c^2.
\]
\end{theorem}
\begin{proof}
Choose the origin at $C$, the positive $x$-axis along $CA$, and the positive $y$-axis parallel to $AB$.

By Axiom~2.1, the coordinate axes represent orthogonal geometric directions. Because $|CA|=a$, vertex $A$ has coordinates $(a,0)$. Because $AB$ is parallel to the $y$-axis with $|AB|=b$, vertex $B$ has coordinates $(a,b)$. The hypotenuse $CB$ has primitive geometric length $c$.

By Axiom~2.2, there exists $R(T)=e^{TJ}\in G$ that maps the ray $CB$ onto the positive $x$-axis, placing $B(a,b)$ at $B'(r,0)$ with $r>0$. We evaluate the terminal axial coordinate $r$ by two channels.

\medskip\noindent\textbf{Channel 1: Dynamical conservation.}
By Lemma~3.1, $Q(x,y)=\lambda(x^2+y^2)$ ($\lambda\ne0$) is conserved along the trajectory from $(a,b)$ to $(r,0)$:
\begin{equation}
Q(a,b)=Q(r,0)\quad\Longrightarrow\quad \lambda(a^2+b^2)=\lambda(r^2+0^2).
\end{equation}
Dividing by $\lambda\ne0$ gives
\begin{equation}\label{eq:qchannel} a^2+b^2=r^2.\end{equation}

\medskip\noindent\textbf{Channel 2: Geometric congruence.}
By Axiom~2.2, rigid rotation preserves primitive geometric length, so $|CB'|=|CB|=c$. By Axiom~2.1, the segment from $(0,0)$ to $(r,0)$ has geometric length $r$. Therefore
\begin{equation}\label{eq:congruence}r=c.\end{equation}
Combining equations~\eqref{eq:qchannel} and~\eqref{eq:congruence} yields
\begin{equation}a^2+b^2=c^2.\end{equation}
\end{proof}

\begin{remark}[Foundational Scope]
This derivation proves the Pythagorean theorem relative to an underlying Euclidean congruence structure and a specified rotational flow, rather than deriving metric geometry from purely algebraic primitives. Its contribution lies in showing that once continuous rotational flow and primitive length congruence are logically decoupled from coordinate metrics, rotational symmetry forces the $x^2+y^2$ quadratic form up to scale, while congruence identifies its terminal axial coordinate with primitive geometric length without assuming the Cartesian distance formula.
\end{remark}

\section{Conclusion}
The coordinate Pythagorean identity follows from the stated Euclidean congruence and rotational-isotropy axioms without assuming an a priori Cartesian diagonal-distance formula. The infinitesimal generator of continuous rotations forces every nonzero invariant symmetric quadratic form to be proportional to $x^2+y^2$, while primitive rigid congruence independently identifies the corresponding terminal axial coordinate with geometric length.

\paragraph{Companion formalization (added in the disclosure revision).}
An independently compiled Lean 4/Mathlib development, including extensions that derive a quarter-turn matrix and generator from additional explicitly stated group and congruence premises, is publicly available at
\url{https://github.com/alexodas/pythagorean-isotropy-lean}.
Compilation verifies the conditional theorems as formalized; it does not establish historical novelty or the minimality of the geometric assumptions.

\begin{thebibliography}{9}
\bibitem{Hall} B. Hall, \emph{Lie Groups, Lie Algebras, and Representations: An Elementary Introduction}, 2nd ed., Graduate Texts in Mathematics, Springer, 2015.
\bibitem{JacksonJohnson} N. Jackson and C. Johnson, "Five or Ten New Proofs of the Pythagorean Theorem," \emph{The American Mathematical Monthly}, vol. 131, no. 9, pp. 739--752, 2024.
\end{thebibliography}

\end{document}
